\documentclass[12pt]{article}
\usepackage{preamble}

\begin{document}

 
%%%%%%%%%%%%%%%%%%%%%%%%%%%%%%%%%%%%%%%%%%%%%%%%%%%%%
%%%%%%%%%%%%%%%%%%%%% Portada %%%%%%%%%%%%%%%%%%%%%%%
%%%%%%%%%%%%%%%%%%%%%%%%%%%%%%%%%%%%%%%%%%%%%%%%%%%%%

\begin{titlepage}
     
    
    % --- LOGOTIPOS (ALINEADOS ARRIBA, IZQUIERDA Y DERECHA) ---
    \noindent
    
\includegraphics[width=0.35\textwidth]{Figuras/ule.jpg}
    \hfill
    
\includegraphics[height=2.5cm]{Figuras/escudo-ingenierias.png}

    \vspace{3cm}

    % --- BLOQUES DE TEXTO CENTRADOS ---
    \begin{center}
        {
        \fontsize{22pt}{26pt}\selectfont\bfseries
        Escuela de Ingenierías \\ Industrial, Informática y Aeroespacial
        }
        
        \vspace{2cm}
        
        {\LARGE\bfseries GRADO EN INGENIERÍA MECÁNICA} \\ 
        \vspace{0.5cm}
        {\LARGE\bfseries INSTALACIONES II}
        
        \vspace{2cm}
        
        {\Large CÁLCULO DE INSTALACIÓN DE GLP}\\[2cm]
    \end{center}
    
    \vfill
    
    % --- AUTORES Y TUTOR (ALINEADOS A LA DERECHA) ---
    \begin{flushright}
        \large
        \begin{tabular}{@{}l@{ }l}
            \textbf{Autores:} & Alberto España Carrete \\
            & Carlos Fernandez Palanca \\
            & Rubén Modino López \\
            & Juan Diego Niño González \\ [5pt]
            \textbf{Tutor:}   & Roberto Getino de la Mano
        \end{tabular}
    \end{flushright}
    
    \vspace{1.5cm}
    
    % --- FECHA (CENTRADA Y CON TAMAÑO ESPECÍFICO DE 22pt) ---
    \begin{center}
        {\fontsize{22pt}{26pt}\selectfont (Mayo, 2026)}
    \end{center}

    \thispagestyle{empty}
\end{titlepage}

%%%% Indices -----------------------------------------------

\hypersetup{linkcolor=black}

\portadacapitulo{Índice General}
\etocsetnexttocdepth{0}
\tableofcontents\thispagestyle{empty}

% En el encabezado se modifican aspectos gen. del indice.

\hypersetup{
    linkcolor=blue,
    citecolor=blue,
    urlcolor=blue
}


%%%%%%%%%%%%%%%%%%%%%%%%%%%%%%%%%%%%%%%%%%%%%%%%%%%%%%%
%%%%%%%%%%%%%%%%%%%%% Contenido %%%%%%%%%%%%%%%%%%%%%%%
%%%%%%%%%%%%%%%%%%%%%%%%%%%%%%%%%%%%%%%%%%%%%%%%%%%%%%%

    \setcounter{section}{1} % Para que la Memoria empiece en el 2
    \portadacapitulo{Memoria}
    \section{Memoria}
    \label{Memoria}

        \subsection{IDENTIFICACIÓN.}
        
        


        \clearpage
        \subsection{OBJETO.}
        ~

        \subsection{ANTECEDENTES.}
        ~

        \subsection{RESUMEN DE CARACTERÍSTICAS.}
            \subsubsection{Titular de la instalación.}
            ~
            \subsubsection{Emplazamiento.}
            ~
            \subsubsection{Localidad.}

            flkjhlfkjs
            ~
            \subsubsection{Tipo de instalación (aérea o enterrada).}
            ~
            \subsubsection{Volumen en \texorpdfstring{m\textsuperscript{3}}{m³} de almacenamiento. Clasificación.}
            ~
            \subsubsection{Instalación receptora. Actividad que utiliza en GLP.}
            ~
            \subsubsection{Potencia térmica total de la instalación en KW.}
            ~
            \subsubsection{Presupuesto total.}
            ~
        
        \subsection{PROMOTOR DE LA INSTALACIÓN.}
        ~
        \subsection{EMPLAZAMIENTO DE LAS INSTALACIONES.}
        ~

        \clearpage
        \subsection{REGLAMENTACIÓN Y NORMAS TÉCNICAS CONSIDERADAS.}
        ~
        \subsection{PLAZO DE EJECUCIÓN.}
        ~
        \subsection{CARACTERÍSTICAS DEL GAS SUMINISTRADO.}
        ~
        \subsection{ACCESOS.}
        ~
        \subsection{DESCRIPCIÓN Y SISTEMA ELEGIDO.}
        
        \subsection{CLASIFICACIÓN Y DISTANCIAS DE SEGURIDAD.}
        ~
        \subsection{CARACTERÍSTICAS DE LOS EQUIPOS.}
            \subsubsection{Depósitos.}
            ~
            \subsubsection{Canalizaciones.}
            ~
            \subsubsection{Boca de carga.}
            ~
            \subsubsection{Equipo de vaporización.}
            ~
            \subsubsection{Equipos de regulación y medida.}
            ~
            \subsubsection{Equipo de trasvase.}
            ~
            \subsubsection{Válvulas de seguridad.}
            ~
        
        \clearpage
        \subsection{INSTALACIÓN ELÉCTRICA.}
        
        
        \subsection{INSTALACIONES DE PROTECCIÓN CONTRAINCENDIOS.}
        ~
        \subsection{PUESTA A TIERRA.}
        ~

    \portadacapitulo{Cálculos Justificados}
    \section{Cálculos justificados}
    \label{Cálculos Justificados}
        \subsection{ÍNDICE DE ANEXOS.}
        ~
        \subsection{CONSUMO Y AUTONOMÍA.}
        ~
        \subsection{VAPORIZACIÓN.}
        ~
        \subsection{VÁLVULAS DE SEGURIDAD.}
        ~
        \subsection{PUNTO MÁXIMO LLENADO: LONGITUD TUBO SONDA.}
        ~
        \subsection{PROTECCIÓN CATÓDICA.}
        ~
        \subsection{PROTECCIÓN CONTRA INCENDIOS.}
        ~


    \portadacapitulo{Presupuesto}
    \section{Presupuesto}
    \label{Presupuesto}
        \subsection{ÍNDICE DE PRESUPUESTO.}
        ~
        \subsection{CUADRO DE PRECIOS Nº 1.}
        ~
        \subsection{CUADRO DE PRECIOS Nº 2.}
        ~
        \subsection{MEDICIONES.}
        ~
        \subsection{PRESUPUESTO.}
        ~

    \portadacapitulo{Planos}
    \label{Planos}
  
        \includepdf[pages=-, fitpaper=true, pagecommand={\thispagestyle{empty}}]{Figuras/situacion-emplazamiento.pdf}
        \includepdf[pages=-, fitpaper=true, pagecommand={\thispagestyle{empty}}]{Figuras/Plano_distancias_seguridad.pdf}
        \includepdf[pages=-, fitpaper=true, pagecommand={\thispagestyle{empty}}]{Figuras/plano_Esquema_instalacion.pdf}
   

    

    %%%% Referencias------------------------

    % \renewcommand{\refname}{Bibliografía} % Cambia el título de "Referencias" a "Bibliografía"
    % \bibliographystyle{IEEEtran}
    % \cleardoublepage
    % \phantomsection
    % \addcontentsline{toc}{section}{Bibliografía}
    % \bibliography{refs}

    %%%% Anexos ----------------------------
    % 2. Iniciamos el entorno de anexos
    %\myappendix{Planos}
        


\end{document}
